\documentclass[10pt,letterpaper]{article}
\usepackage[T1]{fontenc}
\usepackage{lmodern,amsmath,amssymb,amsthm,microtype}
\usepackage[margin=0.9in]{geometry}
\usepackage[hidelinks]{hyperref}
\newtheorem{theorem}{Theorem}
\title{Five colors for rational rotations of a number field}
\author{Aafi Khan Malek\\\small Independent Researcher, Ahmedabad, India\\
\small\href{mailto:aafimalek2032@gmail.com}{aafimalek2032@gmail.com}}
\date{November 2026}
\begin{document}
\maketitle
Let \(K=\mathbb Q(\sqrt{33},\sqrt{-3})\subset\mathbb C\). For a
real rational \(c\in(-1,1)\), let \(\rho=c+i\sqrt{1-c^2}\)
and \(\lambda=2(1-c)\). Write \(v_2(2)=1\).
\begin{theorem}
If \(\rho\notin K\) and \(h=-v_2(\lambda)>0\), the induced
unit-distance graph on \(K\cup\rho K\) has a proper five-coloring.
In particular this holds for
\(\rho=(31+3i\sqrt7)/32\).
\end{theorem}
\paragraph{Scope.}
For \(c=7/8\) the upper bound also follows from the
\href{https://www.erdosproblemaday.com/report/508}{July 2026 working
report on Erd\H{o}s problem 508}, which extends Madore's 11-adic
reduction to the whole coordinate plane over
\(F=\mathbb Q(\sqrt3,\sqrt5,\sqrt{11})\). The example \(c=31/32\)
does not lie in that plane: the imaginary coordinate of \(\rho\)
is \(3\sqrt7/32\), and the prime \(7\) gives a square class
independent of \(3,5,11\). The theorem concerns only the
specified unions, not the full larger coordinate plane,
and gives no new bound for the full Euclidean plane.

\section*{The cross equation}
Embed \(K\) into the unramified quadratic extension
\(\mathbb Q_2(\zeta)\), where \(\zeta^2+\zeta+1=0\), by taking
a \(2\)-adic square root of 33 and sending \(\sqrt{-3}\) to
\(1+2\zeta\). Let \(v(2)=1\) and write \(O=\{x:v(x)\ge0\}\).
The residue field of \(O\) is \(\mathbb F_4\) with kernel \(2O\).
Complex conjugation corresponds to the local involution preserving
valuation, so \(v(x\overline x)=2v(x)\).
The two copies share only zero by \(\rho\notin K\).
If \(x,y\in K\) have a rational cross squared distance, the coefficient
of \(\rho-c\) in its expansion vanishes, since \(1,\rho-c\) are linearly
independent over \(K\). Thus \(z=x\overline y=y\overline x\).
For nonzero endpoints \(t=y/x\) is real and belongs to
\(\mathbb Q(\sqrt{33})\subset\mathbb Q_2\). We have
\[
 |x-\rho y|^2=N(x-y)+\lambda z=N(x)(1-2ct+t^2).
 \tag{1}
\]

\section*{Initial coloring and extension}
First color each integral copy by its \(\mathbb F_4\) residue,
and recolor all nonzero positive-valuation points on one side
with a fifth color. In a native copy unit displacements have
valuation zero. If a cross unit edge initially has equal residues,
then its difference in native coordinates lies in \(2O\).
The first term of (1) has valuation at least two, hence the
product term has valuation zero. This requires \(v(x)+v(y)=h>0\).
If one valuation were zero, they would be different and the
difference would have valuation zero. Both endpoints are
therefore positive-valuation; exactly one was recolored.
The recolored set is independent within its side, and the common
origin was not recolored. This gives proper colors on the integral part.

For a cross unit edge with \(m=\min(v(x),v(y))<0\), no endpoint
is zero. Equal valuations \(m\) are impossible, since the
product term then has valuation \(2m-h\), strictly below
the valuation of the first term, and cannot yield a unit.
Write \(v(x)=m<v(y)\), exchanging the native representatives
if needed. Now the first term has valuation \(2m<0\), so the
second must match it: \(v(y)=m+h\).

Put \(T=2c=2-\lambda\). Then \(v_2(T)=-h\), and
\(1-4/T^2\equiv1\bmod 8\), so the quadratic in (1) has
two roots in \(\mathbb Q_2\). Their valuations are \(-h,h\)
because their sum has valuation \(-h\) and their product is one.
Write \(r_+\) for the root of valuation \(h\). Since
\(v(t)=h\), the distance-one equation gives
\[
 v(t-r_+)=h-2m,\qquad v(y-r_+x)=h-m\ge2. \tag{2}
\]
For any fixed low-side coset \(C=x_0+O\) of valuation \(m<0\),
all higher cross neighbors thus have valuation \(m+h\) and
one common residue class modulo \(2O\). Indeed \(r_+(x-x')\)
has valuation at least \(h\ge1\) for \(x,x'\in C\), and the
errors in (2) have valuation at least two.

Color the cosets at \(m=-1,-2,\ldots\) in the two copies.
The higher cross neighbors of \(C\), if any, have already
received one common color: negative-level coset coloring is
constant on each \(2O\) class, and the integral coloring is
constant there at each fixed valuation. Give the four \(2O\)
classes in \(C\) the other four colors, or any four if no
higher neighbor exists. Native edges stay in one \(O\)-coset
and join distinct residue classes. Every cross edge at negative
valuation is proper by the palette at its lower end. This
proves the theorem.

\paragraph{Effective rule.}
The lower coset representative is obtained by keeping the
finite negative binary digits of the two local coordinates.
A rational approximation \(r'\) to \(r_+\) at precision strictly
larger than \(h\) and at least \(1-m\) gives \(y_*=r'x_0\)
in the same \(2O\) class and of the same valuation as all actual
higher neighbors. Always exclude its color; no neighbor search
is required. The companion exact-rational code checks the rule
on a 171-point, 127-edge fixture for \(c=31/32\), including
edges with negative valuations; this illustrates, but does not
prove, the infinite theorem. The proof is written, not
Lean-formalized. It was developed with AI assistance.

\section*{Optional equality of two integral terminals}
In the integral base coloring one can instead choose the fifth-color
levels individually. A cross unit edge with both representative
valuations positive must have \(a+b=h\), by (1).
For each of the complementary pairs of levels
\((a,h-a)\) on the two sides, \(1\le a<h\), recolor exactly one
entire level. Do not recolor any other level. The union of all
positive levels of either side is independent within that side.

Given two points at Euclidean distance two in the integral union,
if both lie in one side, recolor the opposite side's levels and
retain their common native residue. If the pair is genuinely cross
and both representative valuations are positive, (1) for squared
length four implies \(a+b\ge h+2\). These endpoints are not on
complementary levels, so the shell choices can leave both native
colors unchanged. If one valuation is zero and the other positive,
(1) forces the other to be \(h\); both zero is impossible by the
negative valuation of the product term. Recolor only levels
\(1,\ldots,h-1\) on the positive terminal's side, and give both
terminals the fifth color. The positive terminal cannot neighbor
any positive level in its native side; the zero-valuation one
cannot have a cross unit neighbor of positive valuation less than
\(h\), again by (1). This realizes equality for the specified pair,
without asserting simultaneous equality for all pairs or for cross
pairs at negative valuations.

\end{document}
